\section*{Capítulo \ref{ch:g10:freshwater-resources} --- Os recursos de água doce}

\begin{solution}{exo:g10:freshwater-resources:1}
Um \omterm{def:g10:freshwater-resources:stock}{estoque} é a quantidade de água contida num reservatório num dado
momento, em \unit{km^3} (ou em quilogramas). Um \omterm{def:g10:freshwater-resources:stock}{fluxo} é a quantidade que
passa de um reservatório para outro por unidade de tempo, em \unit{km^3} por
ano. O \omterm{def:g10:freshwater-resources:residence-time}{tempo de residência} é o tempo médio que a água passa num
reservatório, em dias ou em anos; em regime estacionário, ele é o \omterm{def:g10:freshwater-resources:stock}{estoque}
dividido pelo \omterm{def:g10:freshwater-resources:stock}{fluxo}.
\end{solution}

\begin{solution}{exo:g10:freshwater-resources:2}
Cerca de \qty{2.5}{\percent} de toda a água é doce. Da água doce,
\qty{68.7}{\percent} está congelada nas calotas de gelo, nas geleiras e na
neve permanente, e \qty{30.1}{\percent} está no subsolo. Os rios são os que
contêm menos, \num{2120}~\unit{km^3}.
\end{solution}

\begin{solution}{exo:g10:freshwater-resources:3}
A \omterm{def:g10:freshwater-resources:renewable}{água doce renovável} é a água doce renovada a cada ano pela precipitação e
que atravessa uma região como escoamento e como água subterrânea: um volume
por ano, um \omterm{def:g10:freshwater-resources:stock}{fluxo}. A \omterm{def:g10:freshwater-resources:fossil-groundwater}{água subterrânea fóssil} é água que entrou num aquífero
há muito tempo e que a recarga atual não renova: é uma reserva que só pode
ser consumida, um \omterm{def:g10:freshwater-resources:stock}{estoque}.
\end{solution}

\begin{solution}{exo:g10:freshwater-resources:4}
$91\,000/2120 \approx 43$ vezes mais.
\end{solution}

\begin{solution}{exo:g10:freshwater-resources:5}
A atmosfera renova a sua água mais depressa (cerca de \num{9.5} dias), e o
oceano mais devagar entre os reservatórios calculados (cerca de \num{3200}
anos); a água subterrânea profunda é ainda mais antiga.
$3200 \times 365.25/9.5 \approx \num{120000}$ vezes mais tempo.
\end{solution}

\begin{solution}{exo:g10:freshwater-resources:6}
Entrada: $413\,000 + 73\,000 = \num{486000}$~\unit{km^3} por ano, logo
$T = 12\,700/486\,000 \approx \num{0.0261}$ ano $\approx \num{9.5}$ dias.
Saída: $373\,000 + 113\,000 = \num{486000}$~\unit{km^3}, a mesma quantidade:
o ar está em regime estacionário, e o \omterm{def:g10:freshwater-resources:residence-time}{tempo de residência} é o mesmo.
\end{solution}

\begin{solution}{exo:g10:freshwater-resources:7}
$1\,335\,040/413 \approx \num{3230}$ anos. Em regime estacionário, o oceano
recebe tanto quanto perde: $373 + 40 = 413$ mil~\unit{km^3} por ano, de
modo que o \omterm{def:g10:freshwater-resources:residence-time}{tempo de residência} calculado com as entradas é o mesmo.
\end{solution}

\begin{solution}{exo:g10:freshwater-resources:8}
As terras emersas recebem \num{113000}~\unit{km^3} por ano; \num{73000}
voltam ao ar por evaporação e transpiração e \num{40000} escoam para o mar:
$73\,000 + 40\,000 = \num{113000}$. A parte devolvida ao ar é
$73/113 \approx \qty{65}{\percent}$.
\end{solution}

\begin{solution}{exo:g10:freshwater-resources:9}
$273.2 \times 10^6 \times 1233.5 \approx \qty{3.37e11}{m^3}$, isto é,
\qty{337}{km^3}: quase tanto quanto toda a água que o Mar de Aral ainda
continha em 1987.
\end{solution}

\begin{solution}{exo:g10:freshwater-resources:10}
$15.8 \times 0.3048 \approx \qty{4.82}{m}$. Ao longo dos 65 anos de 1950 a
2015, $482/65 \approx \qty{7.4}{cm}$ por ano.
\end{solution}

\begin{solution}{exo:g10:freshwater-resources:11}
$2/0.01 = 200$ anos. O \omterm{def:g10:freshwater-resources:stock}{estoque} agora diminui $0.05 - 0.01 =
\qty{0.04}{km^3}$ por ano, logo dura $2/0.04 = 50$ anos.
\end{solution}

\begin{solution}{exo:g10:freshwater-resources:12}
1960: $1093/71.1 \approx 15$ anos; 2010: $98.1/11.4 \approx 8.6$ anos. A
evaporação retira mais ou menos a mesma espessura de água a cada ano de cada
quilômetro quadrado do lago; à medida que o lago encolheu, a sua área
diminuiu, e com ela o volume total evaporado.
\end{solution}

\begin{solution}{exo:g10:freshwater-resources:13}
(1) A perda líquida é de $1.2 - 0.2 = \qty{1.0}{m}$ por ano sobre
$100 \times 10^6$~\unit{m^2}: $10^8$~\unit{m^3}, ou seja, \qty{0.1}{km^3}
por ano precisa entrar. (2) A entrada cai para \qty{0.05}{km^3}: o lago
perde $\qty{0.05}{km^3} = 5 \times 10^7$~\unit{m^3}, uma camada de
$5 \times 10^7/10^8 = \qty{0.5}{m}$. (3) À medida que o nível desce, as
margens saem da água e a área do lago diminui: menos superfície, menos
evaporação, até que a evaporação e a entrada voltem a se equilibrar.
\end{solution}

\begin{solution}{exo:g10:freshwater-resources:14}
O aquífero contém $100 \times 0.20 = \qty{20}{m}$ de água, isto é,
\qty{20000}{mm}: a \qty{1}{mm} por ano, \num{20000} anos. A água profunda do
aquífero do Saara tem de \num{200000} a um milhão de anos, de dez a
cinquenta vezes mais: ela foi recarregada durante climas passados mais
úmidos.
\end{solution}

\begin{solution}{exo:g10:freshwater-resources:15}
$337/453\,000 \approx \qty{0.00074}{km}$, cerca de \qty{0.74}{m} de água. O
nível da água desceu \qty{4.8}{m}, cerca de seis vezes e meia mais. A água
só ocupa os poros da rocha (e nem toda ela escoa): quando o lençol freático
desce um metro, a água liberada é uma camada muito mais fina que um metro,
aqui cerca de $0.74/4.8 \approx \num{0.15}$~m.
\end{solution}

\begin{solution}{pb:g10:freshwater-resources:1}
\textbf{1.} $1093/68\,000 \approx \qty{0.016}{km}$, cerca de \qty{16}{m}.

\textbf{2.} $1093 \times 10^9 \approx \qty{1.09e12}{m^3}$.

\textbf{3.} \qty{10}{g/L} são \qty{10}{kg/m^3}: $1.093 \times 10^{12} \times 10
\approx \qty{1.1e13}{kg}$, cerca de \num{11} bilhões de toneladas.

\textbf{4.} $1093/71.1 \approx 15$ anos.

\textbf{5.} Nenhum rio sai dele: a sua água só pode sair como vapor (ou se
infiltrando no solo), enquanto o sal fica.

\textbf{6.} $71.1/68\,000 \approx \qty{0.00105}{km}$, cerca de \qty{1.05}{m}.

\textbf{7.} $9.4/68\,000 \approx \qty{0.00014}{km}$, cerca de \qty{0.14}{m}.

\textbf{8.} $71.1 - 9.4 = \qty{61.7}{km^3}$ por ano.

\textbf{9.} $79 + 37 = \qty{116}{km^3}$: $61.7/116 \approx \qty{53}{\percent}$,
cerca da metade.

\textbf{10.} Entradas $1.0 + 6.2 = \qty{7.2}{km^3}$; saída \qty{36.8}{km^3};
variação líquida $7.2 - 36.8 = \qty{-29.6}{km^3}$.

\textbf{11.} $(502.7 - 327.2)/5 = \qty{35.1}{km^3}$ por ano: a mesma ordem
de grandeza que os \qty{29.6}{km^3} do balanço de 1987 (os anos isolados
diferem: a evaporação foi de \qty{47.9}{km^3} em 1984).

\textbf{12.} $36.8/0.00105 \approx \qty{35000}{km^2}$.

\textbf{13.} $35\,000/68\,000 \approx \qty{52}{\percent}$, cerca da metade.

\textbf{14.} O lago tinha encolhido: metade da superfície evapora mais ou
menos metade da água.

\textbf{15.} $345.6/1093 \approx \qty{32}{\percent}$, cerca de um terço.

\textbf{16.} $10 \times 1093/345.6 \approx \qty{31.6}{g/L}$.

\textbf{17.} A salinidade medida é menor: parte do sal tinha saído da água,
cristalizado no leito exposto e sido levado pelo vento como poeira salgada.

\textbf{18.} Cerca de \qty{1.05}{m} por ano, a espessura evaporada.

\textbf{19.} $7.2/35\,000 \approx \qty{0.0002}{km}$, cerca de \qty{0.20}{m}.

\textbf{20.} $1.05 - 0.20 \approx 0.85$ (\num{0.84} com os valores não
arredondados): \textbf{em 1987, o nível do Mar de Aral desceu cerca de
\qty{0.8}{m}, perto de um metro num único ano}.
\end{solution}
